\begin{solapartat}
\punts{0,2} La tercera llei de Kepler estableix que el quadrat del període orbital és
directament proporcional al cub del semieix major de la òrbita. Així doncs, en el cas
d'òrbites circulars s'ha de complir que el quadrat del període orbital és directament
proporcional al cub del radi de la òrbita,

\punts{0,2} Si el planeta descriu un moviment circular uniforme al voltant del Sol, la
seva acceleració centrípeta és: $a_N = \dfrac{v^2}{r}$ o $a_N = \omega^2 r$

I la velocitat es pot expressar en funció del període com: $v = \dfrac{2\pi r}{T}$

Llavors, l'acceleració centrípeta es pot expressar com: $a_N = \left(\dfrac{2\pi}{T}\right)^2 r$
\punts{0,1}

\punts{0,2} I segons la llei de gravitació universal, el mòdul de la força s'expressa com:
\[ F = G\,\frac{m M_\mathrm{Sol}}{r^2} \]

\punts{0,1} La segona llei de Newton estableix que: $\vec{F} = m\vec{a}$, on $m$ és la
massa del planeta.

Per tant, obtenim que: $a = G\,\dfrac{M_\mathrm{Sol}}{r^2}$

\punts{0,2} I com que la força de la gravetat és perpendicular a la trajectòria,
l'acceleració és igual a l'acceleració centrípeta:
\[ G\,\frac{M_\mathrm{Sol}}{r^2} = \left(\frac{2\pi}{T}\right)^2 r \]

I, finalment, si aïllem $T^2$ obtenim la tercera llei de Kepler per a òrbites circulars:
\[ T^2 = \frac{(2\pi)^2}{G M_\mathrm{Sol}} r^3 \ \text{ó}\ \frac{r^3}{T^2} = \frac{G M_\mathrm{Sol}}{(2\pi)^2}. \]
\punts{0,25}
\end{solapartat}

\begin{solapartat}
\punts{1,0} El càlcul es pot fer a partir de l'expressió anterior i fent servir les dades
de qualsevol planeta de la taula:
\[ M_\mathrm{Sol} = \frac{(2\pi)^2}{G}\,\frac{r^3}{T^2} \]
Per exemple, a partir de les dades de l'òrbita de la Terra obtenim:
\[ M_\mathrm{Sol} = \frac{(2\pi)^2}{G}\,\frac{r^3}{T^2} = \frac{(2\pi)^2}{6,67\times 10^{-11}}\,\frac{(149,6\times 10^{9})^3}{(365\times 24\times 3600)^2} = 1,99\times 10^{30}\ \mathrm{kg} \]
\punts{0,25}

També és possible que l'estudiant faci el càlcul per a més d'un planeta i doni el valor
mitjà, aquest procediment també s'ha de considerar correcte.
\end{solapartat}
